% =====================================================================
%  MATH MACROS  —  project-agnostic utility command set
%  File: <project>/style/math-macros.tex
% ---------------------------------------------------------------------
%  Usage in a paper's main.tex (after \documentclass and before
%  \begin{document}):
%
%      % =====================================================================
%  MATH MACROS  —  project-agnostic utility command set
%  File: <project>/style/math-macros.tex
% ---------------------------------------------------------------------
%  Usage in a paper's main.tex (after \documentclass and before
%  \begin{document}):
%
%      % =====================================================================
%  MATH MACROS  —  project-agnostic utility command set
%  File: <project>/style/math-macros.tex
% ---------------------------------------------------------------------
%  Usage in a paper's main.tex (after \documentclass and before
%  \begin{document}):
%
%      % =====================================================================
%  MATH MACROS  —  project-agnostic utility command set
%  File: <project>/style/math-macros.tex
% ---------------------------------------------------------------------
%  Usage in a paper's main.tex (after \documentclass and before
%  \begin{document}):
%
%      \input{style/math-macros}
%
%  Contains only domain-neutral commands that any econ/finance paper
%  in this vault could use. Project-specific symbol vocabulary
%  (\FTPL, \SDF, \HHwealth, ...) lives in notation-canon.tex.
%
%  Self-contained: assumes only what workingpaper.cls already loads
%  (amsmath, amsthm, bm, newtxmath). No new \usepackage calls.
% =====================================================================


% ---------------------------------------------------------------------
%  Probability and expectation operators
%
%  Convention: blackboard-bold E and P for measure-theoretic objects;
%  upright Var/Cov for moment statistics.
% ---------------------------------------------------------------------
\newcommand{\E}{\mathbb{E}}
\newcommand{\Prob}{\mathbb{P}}
\newcommand{\Var}{\mathrm{Var}}
\newcommand{\Cov}{\mathrm{Cov}}
\newcommand{\Corr}{\mathrm{Corr}}

% Conditional shortcuts: \Et[X], \Vart[X], \Covt[X,Y].
% (Subscript-as-token, so \Et^{\F_s} also works.)
\newcommand{\Et}{\E_{t}}
\newcommand{\Vart}{\Var_{t}}
\newcommand{\Covt}{\Cov_{t}}


% ---------------------------------------------------------------------
%  Number sets
% ---------------------------------------------------------------------
\newcommand{\R}{\mathbb{R}}
\newcommand{\N}{\mathbb{N}}
\newcommand{\Z}{\mathbb{Z}}
\newcommand{\Q}{\mathbb{Q}}
\newcommand{\Cset}{\mathbb{C}}        % \C is taken by some packages; alias
\newcommand{\Rplus}{\R_{+}}
\newcommand{\Rpp}{\R_{++}}            % strictly positive reals


% ---------------------------------------------------------------------
%  Distributions and convergence
% ---------------------------------------------------------------------
\newcommand{\Normal}{\mathcal{N}}                     % Gaussian
\newcommand{\Unif}{\mathcal{U}}                       % uniform
\newcommand{\iidsim}{\overset{\text{i.i.d.}}{\sim}}   % X_i \iidsim \Normal(0,1)
\newcommand{\todist}{\overset{d}{\longrightarrow}}    % convergence in distribution
\newcommand{\toprob}{\overset{p}{\longrightarrow}}    % convergence in probability
\newcommand{\toas}{\overset{a.s.}{\longrightarrow}}   % almost-sure convergence


% ---------------------------------------------------------------------
%  Differential calculus
%
%  \dd produces an upright "d" with thin-space spacing for integrals:
%      \int_0^1 f(x) \dd x
%
%  \pd / \pdtwo wrap the partial-derivative fraction:
%      \pd{f}{x}        ->  df/dx (partial)
%      \pdtwo{f}{x}     ->  d^2 f / dx^2  (partial)
%      \pdmix{f}{x}{y}  ->  d^2 f / dx dy  (mixed partial)
% ---------------------------------------------------------------------
\newcommand{\dd}{\mathop{}\!\mathrm{d}}
\newcommand{\pd}[2]{\frac{\partial #1}{\partial #2}}
\newcommand{\pdtwo}[2]{\frac{\partial^{2} #1}{\partial #2^{2}}}
\newcommand{\pdmix}[3]{\frac{\partial^{2} #1}{\partial #2 \, \partial #3}}


% ---------------------------------------------------------------------
%  Delimiters
%
%  Plain (non-auto-sized) versions. For auto-sizing in display math,
%  use \left\lvert ... \right\rvert directly, or load mathtools in
%  the paper's preamble and switch to \DeclarePairedDelimiter.
% ---------------------------------------------------------------------
\newcommand{\abs}[1]{\lvert #1 \rvert}
\newcommand{\norm}[1]{\lVert #1 \rVert}
\newcommand{\inner}[2]{\langle #1, #2 \rangle}
\newcommand{\set}[1]{\{ #1 \}}                        % set-builder shorthand
\newcommand{\given}{\;\middle|\;}                     % conditional bar in P(A \given B)


% ---------------------------------------------------------------------
%  Indicator function and common matrix ops
%
%  \argmax, \argmin, \plim are already declared in workingpaper.cls;
%  do not duplicate them here.
% ---------------------------------------------------------------------
\newcommand{\indic}{\mathbf{1}}                       % \indic\{X > 0\}
\DeclareMathOperator{\sign}{sign}
\DeclareMathOperator{\diag}{diag}
\DeclareMathOperator{\tr}{tr}
\DeclareMathOperator{\rk}{rank}
\DeclareMathOperator{\spn}{span}
\DeclareMathOperator{\Span}{span}                     % alias
\DeclareMathOperator{\proj}{proj}


% ---------------------------------------------------------------------
%  Vector and matrix typography (uses bm, loaded by workingpaper.cls)
%
%  Convention: \vect for column vectors, \mat for matrices. Both apply
%  bold-italic via \bm so that letters retain math-italic shape.
%  For upright bold (e.g. canonical R^n basis vectors), use \mathbf.
% ---------------------------------------------------------------------
\newcommand{\vect}[1]{\bm{#1}}
\newcommand{\mat}[1]{\bm{#1}}
\newcommand{\transpose}{^{\top}}                      % X\transpose -> X^T
\newcommand{\inv}{^{-1}}                              % X\inv      -> X^{-1}


% ---------------------------------------------------------------------
%  Spacing helpers
% ---------------------------------------------------------------------
\newcommand{\eqsep}{\;}                               % medium space inside equations
\newcommand{\where}{\quad\text{where }}               % "where" with spacing
\newcommand{\suchthat}{\;:\;}                         % shortcut for set-builder colon


\endinput

%
%  Contains only domain-neutral commands that any econ/finance paper
%  in this vault could use. Project-specific symbol vocabulary
%  (\FTPL, \SDF, \HHwealth, ...) lives in notation-canon.tex.
%
%  Self-contained: assumes only what workingpaper.cls already loads
%  (amsmath, amsthm, bm, newtxmath). No new \usepackage calls.
% =====================================================================


% ---------------------------------------------------------------------
%  Probability and expectation operators
%
%  Convention: blackboard-bold E and P for measure-theoretic objects;
%  upright Var/Cov for moment statistics.
% ---------------------------------------------------------------------
\newcommand{\E}{\mathbb{E}}
\newcommand{\Prob}{\mathbb{P}}
\newcommand{\Var}{\mathrm{Var}}
\newcommand{\Cov}{\mathrm{Cov}}
\newcommand{\Corr}{\mathrm{Corr}}

% Conditional shortcuts: \Et[X], \Vart[X], \Covt[X,Y].
% (Subscript-as-token, so \Et^{\F_s} also works.)
\newcommand{\Et}{\E_{t}}
\newcommand{\Vart}{\Var_{t}}
\newcommand{\Covt}{\Cov_{t}}


% ---------------------------------------------------------------------
%  Number sets
% ---------------------------------------------------------------------
\newcommand{\R}{\mathbb{R}}
\newcommand{\N}{\mathbb{N}}
\newcommand{\Z}{\mathbb{Z}}
\newcommand{\Q}{\mathbb{Q}}
\newcommand{\Cset}{\mathbb{C}}        % \C is taken by some packages; alias
\newcommand{\Rplus}{\R_{+}}
\newcommand{\Rpp}{\R_{++}}            % strictly positive reals


% ---------------------------------------------------------------------
%  Distributions and convergence
% ---------------------------------------------------------------------
\newcommand{\Normal}{\mathcal{N}}                     % Gaussian
\newcommand{\Unif}{\mathcal{U}}                       % uniform
\newcommand{\iidsim}{\overset{\text{i.i.d.}}{\sim}}   % X_i \iidsim \Normal(0,1)
\newcommand{\todist}{\overset{d}{\longrightarrow}}    % convergence in distribution
\newcommand{\toprob}{\overset{p}{\longrightarrow}}    % convergence in probability
\newcommand{\toas}{\overset{a.s.}{\longrightarrow}}   % almost-sure convergence


% ---------------------------------------------------------------------
%  Differential calculus
%
%  \dd produces an upright "d" with thin-space spacing for integrals:
%      \int_0^1 f(x) \dd x
%
%  \pd / \pdtwo wrap the partial-derivative fraction:
%      \pd{f}{x}        ->  df/dx (partial)
%      \pdtwo{f}{x}     ->  d^2 f / dx^2  (partial)
%      \pdmix{f}{x}{y}  ->  d^2 f / dx dy  (mixed partial)
% ---------------------------------------------------------------------
\newcommand{\dd}{\mathop{}\!\mathrm{d}}
\newcommand{\pd}[2]{\frac{\partial #1}{\partial #2}}
\newcommand{\pdtwo}[2]{\frac{\partial^{2} #1}{\partial #2^{2}}}
\newcommand{\pdmix}[3]{\frac{\partial^{2} #1}{\partial #2 \, \partial #3}}


% ---------------------------------------------------------------------
%  Delimiters
%
%  Plain (non-auto-sized) versions. For auto-sizing in display math,
%  use \left\lvert ... \right\rvert directly, or load mathtools in
%  the paper's preamble and switch to \DeclarePairedDelimiter.
% ---------------------------------------------------------------------
\newcommand{\abs}[1]{\lvert #1 \rvert}
\newcommand{\norm}[1]{\lVert #1 \rVert}
\newcommand{\inner}[2]{\langle #1, #2 \rangle}
\newcommand{\set}[1]{\{ #1 \}}                        % set-builder shorthand
\newcommand{\given}{\;\middle|\;}                     % conditional bar in P(A \given B)


% ---------------------------------------------------------------------
%  Indicator function and common matrix ops
%
%  \argmax, \argmin, \plim are already declared in workingpaper.cls;
%  do not duplicate them here.
% ---------------------------------------------------------------------
\newcommand{\indic}{\mathbf{1}}                       % \indic\{X > 0\}
\DeclareMathOperator{\sign}{sign}
\DeclareMathOperator{\diag}{diag}
\DeclareMathOperator{\tr}{tr}
\DeclareMathOperator{\rk}{rank}
\DeclareMathOperator{\spn}{span}
\DeclareMathOperator{\Span}{span}                     % alias
\DeclareMathOperator{\proj}{proj}


% ---------------------------------------------------------------------
%  Vector and matrix typography (uses bm, loaded by workingpaper.cls)
%
%  Convention: \vect for column vectors, \mat for matrices. Both apply
%  bold-italic via \bm so that letters retain math-italic shape.
%  For upright bold (e.g. canonical R^n basis vectors), use \mathbf.
% ---------------------------------------------------------------------
\newcommand{\vect}[1]{\bm{#1}}
\newcommand{\mat}[1]{\bm{#1}}
\newcommand{\transpose}{^{\top}}                      % X\transpose -> X^T
\newcommand{\inv}{^{-1}}                              % X\inv      -> X^{-1}


% ---------------------------------------------------------------------
%  Spacing helpers
% ---------------------------------------------------------------------
\newcommand{\eqsep}{\;}                               % medium space inside equations
\newcommand{\where}{\quad\text{where }}               % "where" with spacing
\newcommand{\suchthat}{\;:\;}                         % shortcut for set-builder colon


\endinput

%
%  Contains only domain-neutral commands that any econ/finance paper
%  in this vault could use. Project-specific symbol vocabulary
%  (\FTPL, \SDF, \HHwealth, ...) lives in notation-canon.tex.
%
%  Self-contained: assumes only what workingpaper.cls already loads
%  (amsmath, amsthm, bm, newtxmath). No new \usepackage calls.
% =====================================================================


% ---------------------------------------------------------------------
%  Probability and expectation operators
%
%  Convention: blackboard-bold E and P for measure-theoretic objects;
%  upright Var/Cov for moment statistics.
% ---------------------------------------------------------------------
\newcommand{\E}{\mathbb{E}}
\newcommand{\Prob}{\mathbb{P}}
\newcommand{\Var}{\mathrm{Var}}
\newcommand{\Cov}{\mathrm{Cov}}
\newcommand{\Corr}{\mathrm{Corr}}

% Conditional shortcuts: \Et[X], \Vart[X], \Covt[X,Y].
% (Subscript-as-token, so \Et^{\F_s} also works.)
\newcommand{\Et}{\E_{t}}
\newcommand{\Vart}{\Var_{t}}
\newcommand{\Covt}{\Cov_{t}}


% ---------------------------------------------------------------------
%  Number sets
% ---------------------------------------------------------------------
\newcommand{\R}{\mathbb{R}}
\newcommand{\N}{\mathbb{N}}
\newcommand{\Z}{\mathbb{Z}}
\newcommand{\Q}{\mathbb{Q}}
\newcommand{\Cset}{\mathbb{C}}        % \C is taken by some packages; alias
\newcommand{\Rplus}{\R_{+}}
\newcommand{\Rpp}{\R_{++}}            % strictly positive reals


% ---------------------------------------------------------------------
%  Distributions and convergence
% ---------------------------------------------------------------------
\newcommand{\Normal}{\mathcal{N}}                     % Gaussian
\newcommand{\Unif}{\mathcal{U}}                       % uniform
\newcommand{\iidsim}{\overset{\text{i.i.d.}}{\sim}}   % X_i \iidsim \Normal(0,1)
\newcommand{\todist}{\overset{d}{\longrightarrow}}    % convergence in distribution
\newcommand{\toprob}{\overset{p}{\longrightarrow}}    % convergence in probability
\newcommand{\toas}{\overset{a.s.}{\longrightarrow}}   % almost-sure convergence


% ---------------------------------------------------------------------
%  Differential calculus
%
%  \dd produces an upright "d" with thin-space spacing for integrals:
%      \int_0^1 f(x) \dd x
%
%  \pd / \pdtwo wrap the partial-derivative fraction:
%      \pd{f}{x}        ->  df/dx (partial)
%      \pdtwo{f}{x}     ->  d^2 f / dx^2  (partial)
%      \pdmix{f}{x}{y}  ->  d^2 f / dx dy  (mixed partial)
% ---------------------------------------------------------------------
\newcommand{\dd}{\mathop{}\!\mathrm{d}}
\newcommand{\pd}[2]{\frac{\partial #1}{\partial #2}}
\newcommand{\pdtwo}[2]{\frac{\partial^{2} #1}{\partial #2^{2}}}
\newcommand{\pdmix}[3]{\frac{\partial^{2} #1}{\partial #2 \, \partial #3}}


% ---------------------------------------------------------------------
%  Delimiters
%
%  Plain (non-auto-sized) versions. For auto-sizing in display math,
%  use \left\lvert ... \right\rvert directly, or load mathtools in
%  the paper's preamble and switch to \DeclarePairedDelimiter.
% ---------------------------------------------------------------------
\newcommand{\abs}[1]{\lvert #1 \rvert}
\newcommand{\norm}[1]{\lVert #1 \rVert}
\newcommand{\inner}[2]{\langle #1, #2 \rangle}
\newcommand{\set}[1]{\{ #1 \}}                        % set-builder shorthand
\newcommand{\given}{\;\middle|\;}                     % conditional bar in P(A \given B)


% ---------------------------------------------------------------------
%  Indicator function and common matrix ops
%
%  \argmax, \argmin, \plim are already declared in workingpaper.cls;
%  do not duplicate them here.
% ---------------------------------------------------------------------
\newcommand{\indic}{\mathbf{1}}                       % \indic\{X > 0\}
\DeclareMathOperator{\sign}{sign}
\DeclareMathOperator{\diag}{diag}
\DeclareMathOperator{\tr}{tr}
\DeclareMathOperator{\rk}{rank}
\DeclareMathOperator{\spn}{span}
\DeclareMathOperator{\Span}{span}                     % alias
\DeclareMathOperator{\proj}{proj}


% ---------------------------------------------------------------------
%  Vector and matrix typography (uses bm, loaded by workingpaper.cls)
%
%  Convention: \vect for column vectors, \mat for matrices. Both apply
%  bold-italic via \bm so that letters retain math-italic shape.
%  For upright bold (e.g. canonical R^n basis vectors), use \mathbf.
% ---------------------------------------------------------------------
\newcommand{\vect}[1]{\bm{#1}}
\newcommand{\mat}[1]{\bm{#1}}
\newcommand{\transpose}{^{\top}}                      % X\transpose -> X^T
\newcommand{\inv}{^{-1}}                              % X\inv      -> X^{-1}


% ---------------------------------------------------------------------
%  Spacing helpers
% ---------------------------------------------------------------------
\newcommand{\eqsep}{\;}                               % medium space inside equations
\newcommand{\where}{\quad\text{where }}               % "where" with spacing
\newcommand{\suchthat}{\;:\;}                         % shortcut for set-builder colon


\endinput

%
%  Contains only domain-neutral commands that any econ/finance paper
%  in this vault could use. Project-specific symbol vocabulary
%  (\FTPL, \SDF, \HHwealth, ...) lives in notation-canon.tex.
%
%  Self-contained: assumes only what workingpaper.cls already loads
%  (amsmath, amsthm, bm, newtxmath). No new \usepackage calls.
% =====================================================================


% ---------------------------------------------------------------------
%  Probability and expectation operators
%
%  Convention: blackboard-bold E and P for measure-theoretic objects;
%  upright Var/Cov for moment statistics.
% ---------------------------------------------------------------------
\newcommand{\E}{\mathbb{E}}
\newcommand{\Prob}{\mathbb{P}}
\newcommand{\Var}{\mathrm{Var}}
\newcommand{\Cov}{\mathrm{Cov}}
\newcommand{\Corr}{\mathrm{Corr}}

% Conditional shortcuts: \Et[X], \Vart[X], \Covt[X,Y].
% (Subscript-as-token, so \Et^{\F_s} also works.)
\newcommand{\Et}{\E_{t}}
\newcommand{\Vart}{\Var_{t}}
\newcommand{\Covt}{\Cov_{t}}


% ---------------------------------------------------------------------
%  Number sets
% ---------------------------------------------------------------------
\newcommand{\R}{\mathbb{R}}
\newcommand{\N}{\mathbb{N}}
\newcommand{\Z}{\mathbb{Z}}
\newcommand{\Q}{\mathbb{Q}}
\newcommand{\Cset}{\mathbb{C}}        % \C is taken by some packages; alias
\newcommand{\Rplus}{\R_{+}}
\newcommand{\Rpp}{\R_{++}}            % strictly positive reals


% ---------------------------------------------------------------------
%  Distributions and convergence
% ---------------------------------------------------------------------
\newcommand{\Normal}{\mathcal{N}}                     % Gaussian
\newcommand{\Unif}{\mathcal{U}}                       % uniform
\newcommand{\iidsim}{\overset{\text{i.i.d.}}{\sim}}   % X_i \iidsim \Normal(0,1)
\newcommand{\todist}{\overset{d}{\longrightarrow}}    % convergence in distribution
\newcommand{\toprob}{\overset{p}{\longrightarrow}}    % convergence in probability
\newcommand{\toas}{\overset{a.s.}{\longrightarrow}}   % almost-sure convergence


% ---------------------------------------------------------------------
%  Differential calculus
%
%  \dd produces an upright "d" with thin-space spacing for integrals:
%      \int_0^1 f(x) \dd x
%
%  \pd / \pdtwo wrap the partial-derivative fraction:
%      \pd{f}{x}        ->  df/dx (partial)
%      \pdtwo{f}{x}     ->  d^2 f / dx^2  (partial)
%      \pdmix{f}{x}{y}  ->  d^2 f / dx dy  (mixed partial)
% ---------------------------------------------------------------------
\newcommand{\dd}{\mathop{}\!\mathrm{d}}
\newcommand{\pd}[2]{\frac{\partial #1}{\partial #2}}
\newcommand{\pdtwo}[2]{\frac{\partial^{2} #1}{\partial #2^{2}}}
\newcommand{\pdmix}[3]{\frac{\partial^{2} #1}{\partial #2 \, \partial #3}}


% ---------------------------------------------------------------------
%  Delimiters
%
%  Plain (non-auto-sized) versions. For auto-sizing in display math,
%  use \left\lvert ... \right\rvert directly, or load mathtools in
%  the paper's preamble and switch to \DeclarePairedDelimiter.
% ---------------------------------------------------------------------
\newcommand{\abs}[1]{\lvert #1 \rvert}
\newcommand{\norm}[1]{\lVert #1 \rVert}
\newcommand{\inner}[2]{\langle #1, #2 \rangle}
\newcommand{\set}[1]{\{ #1 \}}                        % set-builder shorthand
\newcommand{\given}{\;\middle|\;}                     % conditional bar in P(A \given B)


% ---------------------------------------------------------------------
%  Indicator function and common matrix ops
%
%  \argmax, \argmin, \plim are already declared in workingpaper.cls;
%  do not duplicate them here.
% ---------------------------------------------------------------------
\newcommand{\indic}{\mathbf{1}}                       % \indic\{X > 0\}
\DeclareMathOperator{\sign}{sign}
\DeclareMathOperator{\diag}{diag}
\DeclareMathOperator{\tr}{tr}
\DeclareMathOperator{\rk}{rank}
\DeclareMathOperator{\spn}{span}
\DeclareMathOperator{\Span}{span}                     % alias
\DeclareMathOperator{\proj}{proj}


% ---------------------------------------------------------------------
%  Vector and matrix typography (uses bm, loaded by workingpaper.cls)
%
%  Convention: \vect for column vectors, \mat for matrices. Both apply
%  bold-italic via \bm so that letters retain math-italic shape.
%  For upright bold (e.g. canonical R^n basis vectors), use \mathbf.
% ---------------------------------------------------------------------
\newcommand{\vect}[1]{\bm{#1}}
\newcommand{\mat}[1]{\bm{#1}}
\newcommand{\transpose}{^{\top}}                      % X\transpose -> X^T
\newcommand{\inv}{^{-1}}                              % X\inv      -> X^{-1}


% ---------------------------------------------------------------------
%  Spacing helpers
% ---------------------------------------------------------------------
\newcommand{\eqsep}{\;}                               % medium space inside equations
\newcommand{\where}{\quad\text{where }}               % "where" with spacing
\newcommand{\suchthat}{\;:\;}                         % shortcut for set-builder colon


\endinput
